\scanpage{13}
% archive sheet number: 12

Dans ce cas favorable
\[
 \sum_i(-1)^i\frac{i+\rho}{2}b_i
 =\frac{\rho}{2}\chi(E)+\frac12
 \underset{\substack{\rule{0.4pt}{4pt}\rule{6em}{0.4pt}\rule{0.4pt}{4pt}\\ |}}
 {\sum_{i=0}^{2n}(-1)^i i b_i}
\]
\[
 \sum_{i=0}^{n-1}\bigl[(-1)^i i b_i+
 (\deleted{\unsure{2|-2}}-1)^{2n-i}(2n-i)b_{2n-i}\bigr]+(-1)^n n b_n
\]
comme $b_i=b_{2n-i}$
\[
 2n\sum_{i=0}^{n-1}(-1)^i b_i+(-1)^n n b_n
 =\deleted{\unsure{2|+}}n\sum_{i=0}^{2n}(-1)^i b_i=n\chi(E)
\]
\[
 \sum_i(-1)^i\frac{i+\rho}{2}b_i
 =\deleted{n+\rho}\frac{n+\rho}{2}\chi(E) \qquad \text{donc}
\]
\[
 \delta(E)=\deleted{\unsure{\delta|\varepsilon}}\unsure{\delta|\varepsilon}(E)\,
 p^{\frac{n+\rho}{2}\chi(E)}
 \qquad \varepsilon(E)=\pm1
\]

Par suite l'équation fonctionnelle devient dans ce cas
\begingroup\csname tagsleft@true\endcsname\[
 L_{DE}(t)=\varepsilon(E)\left(-p^{\frac{n+\rho}{2}}\frac{1}{t}\right)^{\chi(E)}
 L_E(t^{-1}) \tag*{(2)}
\]\endgroup

Ou encore, si on admet que
\[
 \check E\deleted{\unsure{\otimes E|\simeq E}}\simeq E(\rho)
 \qquad \text{globalement sur }X,\ \text{\deleted{\unsure{d'où|donc}}}
\]
[il suffit que ce soit vrai sur le corps des \unsure{fonctions|fractions} de $X$, donc c'est \unsure{sans doute} tjrs vrai]
\[
 \deleted{\unsure{L_E|L_{\check E}}}D(E)\simeq
 \deleted{\check E\otimes\unsure{\mathbb Q_\ell(n)|\mathbb Q_\ell}}
 \check E(n)\underset{\substack{\diagup\\\text{shift}}}{[2n]}
 \simeq E(\rho+n)\underset{\substack{\diagup\\\text{shift.}}}{[2n]}
\]
et par suite
\[
 \deleted{L_{DE}\simeq L_{E(\rho+n)}=
 L_{\unsure{\mathbb R|D|R}f_* E(\rho+n)}}
\]
\[
 \mathbb R f_{\unsure{{}|*}}(D(E))\simeq
 \mathbb R f_*(E)(\rho+\deleted{\unsure{2|3}}n)[2n]\qquad \text{d'où}
\]
\[
 L_{DE}\unsure{(t)|{}}=L_{\mathbb R f_*D(E)}\unsure{(t)|{}}
 \simeq L_{\mathbb R f_*(E)}(p^{-(\rho+n)}t)
\]
\deleted{et il \inserted{\unsure{s'en|s'e}} \unsure{suit|vient}}

